% Bölüm 20 — taslak plan. Tam metin yazıldığında bu dosyanın yerini alır.
\chapter{Karar Verilebilirlik ve Otomasyon}\label{bol:otomasyon}

\begin{bolumhedefleri}
\begin{itemize}
  \item Karar verilebilirlik ve Entscheidungsproblem'i
  \item Decidable sınıfını ve decide taktiğini
  \item SAT ve SMT çözücülerin mantıksal temellerini
  \item Lean'deki otomasyon araçlarını
\end{itemize}
\end{bolumhedefleri}

\begin{taslakplan}
\textbf{Planlanan kesimler}
\begin{enumerate}[label=\thechapter.\arabic*, leftmargin=2.4em]
  \item Entscheidungsproblem ve Turing
  \item Decidable ve decide
  \item Önermeler mantığı için karar yöntemi
  \item SAT, SMT ve Lean bağlantıları
  \item simp, omega, aesop, polyrith
\end{enumerate}

\textbf{Planlanan Lean örnekleri}
\begin{itemize}
  \item \lc{Totoloji denetleyici yazmak ve doğruluğunu ispatlamak}
\end{itemize}

\textbf{Temel kaynaklar:} \citet{turing1936,harrison2009,demoura2008z3}.
\end{taslakplan}

